% status report 5
\documentclass[12pt]{article}
\usepackage{graphicx}
\usepackage{fancyhdr}
\usepackage{geometry}
\usepackage{hyperref}
\usepackage{xcolor}
\usepackage{url}
\geometry{margin=1in}

\title{CSE 123A - Week 5 Status}
\author{Group 8}
\date{}

\begin{document}
\pagestyle{fancy}
\maketitle

\fancyfoot{}
\fancyhead{}
\fancyfoot[R]{\thepage}

 % whats been completed 
\section*{What we Did}


 
% whats in progress
\section*{What we are Working On}


% what next on teh agenda 
\section*{What's next}
We plan to print the newly redesigned base plates, which has a flat bottom to provide more accurate and stable weight 
readings. We are also increasing the plate thickness to handle heavier loads and expanding the plate dimensions to accommodate 
a wider variety of container sizes.

% also have indivudal secitons of what each persion did and is working on

\section*{Individual Contributions}

\noindent\textbf{Dev Chodavadiya}:
This week, I redesigned the base plate to increase its thickness and revised the screw hole layout to a racetrack pattern, 
ensuring the bottom surface sits flat. I also created engineering drawings for the base plates and worked on developing a 
demo script for our end-of-quarter initial prototype presentation.\\

\noindent\textbf{Logan Bossuwe}:


\noindent\textbf{Pieter Nauwelaerts}:
Significant updates to diagrams for the initial prototype presentation. Changes include noting that we will be showing the smart base reporting data to a laptop serial connection rather than through a WiFi to server connection. We will also be demonstrating a web app rather than a working phone app. 

I also began the process of creating test programs for the current sensor setup so we can test the code written and determine if any fixes or possible corner cases were left out. This also includes limit testing how fast the data can be updated, as well as how it handles repeated changes to the state. 

\noindent\textbf{Reid Graham}:
This week, I restructured the design document to better match the outline in Canvas and to include important information on the architecture of our product.

 I also updated the architecture diagram after our discussion on communication between the microcontroller and web server. It now shows that the MCU will only send data, the weight, and calculations and calibration will happen on the webapp side. \\ 

\noindent\textbf{Sam Lai}:
This week, I worked on creating unit tests for the web app, specifically for the water level calculation and calibration actions. 
The unit tests cover several cases, such as when there is no weight data available, whether the water level is correctly calculated 
based on the weight data, and whether the calibration action correctly updates the calibration data. I am currently working on reimplementing 
the notifications using Push API instead of Firebase Cloud Messaging. Push API has all the features of Firebase Cloud Messaging, but 
it is able to send push notifications on IOS devices, which is a major advantage over Firebase Cloud Messaging.

\end{document}
